\chapter{Children, Families and Decisions About Risk}
\label{ch:family}

\section{The hardest decisions}

Some of the most consequential decisions courts make concern children. Should a child be removed from parents who may harm her? Should a care order be made, or an adoption order that permanently ends the legal relationship with the birth family? With which parent should a child live after separation, and how much contact should the other have? These decisions are made under uncertainty about past events, such as how an injury occurred, and about future risks, such as whether a parent will harm a child again. The losses from error are severe on both sides: a child left with an abusive parent may be injured or killed; a child wrongly removed loses her family. Family law is therefore a domain where the decision theory of this book applies with particular force, and where courts have, to a striking degree, articulated their reasoning in probabilistic terms, although that language appears in the tests for future harm and for the pool of possible perpetrators rather than in the finding of past facts, which remains binary.\index{family law}\index{care proceedings}

\section{The threshold criteria}

Under section 31 of the Children Act 1989, which governs England and Wales, a court may make a care or supervision order only if it is satisfied that the child is suffering, or is likely to suffer, significant harm, attributable to the care given to the child, or likely to be given if the order were not made, not being what it would be reasonable to expect a parent to give, or to the child being beyond parental control. In \emph{Re H (Minors)}, discussed in Chapter~\ref{ch:standards}, the House of Lords held that likely means a real possibility\index{real possibility}, a possibility that cannot sensibly be ignored, rather than more probable than not, but that any finding of likelihood must be based on facts proved on the balance of probabilities, not on suspicion \citep{reh1996}.\index{threshold criteria}

The structure has two levels. Past facts must be proved on the balance of probabilities: did this injury occur, and was it inflicted? Future harm need only be a real possibility given the proved facts. In decision terms, the first level is a threshold on a posterior about the past; the second is a lower threshold on a predictive probability about the future. The lower threshold for future harm reflects the severity of the harm the law seeks to prevent: if the loss from leaving a child at risk is very large relative to the loss from intervening, intervention is justified at a modest probability of harm, as is consistent with the analysis of Chapter~\ref{ch:decisions}. The requirement that the predictive probability rest on proved facts, not suspicion, protects families from intervention on speculation. It is not, however, a Bayesian rule. The court conditions the prediction only on the facts it has found proved on the balance of probabilities, and a past fact found at a probability of $0.4$ is treated as though it had probability zero, whereas a Bayesian predictive probability would average over the unproven facts in proportion to their probabilities. The rule discards probability, a point pressed by Lord Browne-Wilkinson in his dissent in \emph{Re H}, and it can be defended only on grounds outside the calculus, such as the danger of intervention on suspicion and the need for a fixed factual footing for the order. The threshold on the predictive probability is thus a threshold on a probability conditional on the proved facts alone.

\section{Uncertain perpetrators}

When a child has been injured and several adults had care of the child, the court may be unable to identify which of them caused the injury. In \emph{Lancashire County Council v B} the House of Lords held that the threshold could be satisfied where the harm was attributable to the care given by one of the carers, even if it could not be determined which \citep{lancashire2000}. In \emph{Re S-B (Children)} the Supreme Court held that a court should identify the perpetrator if it can on the balance of probabilities, but that if it cannot, it should identify the pool of possible perpetrators, those of whom there is a likelihood or real possibility that they were the perpetrator \citep{resb2009}.\index{uncertain perpetrator}

This is the island problem of Chapter~\ref{ch:classic} in a family court. The pool of possible perpetrators is the set of alternatives, each with a posterior probability. If one carer's posterior exceeds one half, that carer is found to be the perpetrator. If none does, but the posterior is spread over two or three people, each carer for whom there is a real possibility is in the pool. The analogy assumes that exactly one carer is the perpetrator, so that the alternatives are mutually exclusive. The subsequent decision, whether the child can safely return to one of them, depends on that person's posterior probability of having caused the injury and on the risk the person would pose in future, combined with the losses from each decision. A carer with a posterior of forty per cent is neither found to be the perpetrator nor excluded; the court must decide how to protect the child given that uncertainty.

\section{Expert evidence of abuse}

Many care cases turn on medical evidence that an injury was inflicted rather than accidental. Some patterns of injury, such as certain fractures in infants or the combination of findings once called shaken baby syndrome, have been treated as strongly indicative of abuse. The reliability of these inferences has been contested, and several criminal convictions based on them have been overturned. The Bayesian analysis is the one of Chapter~\ref{ch:evidence}: a finding is evidence of abuse to the extent that it is more probable in abused than in non abused children with similar presentations. The relevant comparison group matters, as does the possibility of rare medical conditions that mimic abuse. Experts who state that a finding is diagnostic of abuse are claiming a likelihood ratio close to infinity, which few medical findings can support. The Court of Appeal in \emph{Re U; Re B} held that the civil standard applies in care proceedings and warned judges to be on guard against the over-dogmatic expert \citep{reu2004}. That is a caution about how expert opinion is weighed, not a rule requiring experts to explain its basis, but it points the same way as the Royal Statistical Society's report on investigating suspected medical misconduct, to which Chapter~\ref{ch:clusters} returns.

\section{Proportionality and adoption}

In \emph{Re B (Care Proceedings: Appeal)} the Supreme Court held in 2013 that a care order with a plan for adoption, which severs the relationship between child and birth parents, is a very extreme thing, to be made only where nothing else will do \citep{reb2013}. The Court of Appeal in \emph{Re B-S (Children)} required courts to undertake a global, holistic evaluation of each realistic option, weighing the advantages and disadvantages of each, rather than a linear process that rejects options one by one \citep{rebs2013}.\index{adoption}

\emph{Re B-S} is a call for decision analysis. Each realistic option, return to a parent with support, placement with relatives, long term foster care, adoption, has an uncertain outcome for the child. A holistic evaluation compares their expected outcomes, taking account of the probabilities of success and failure for each and the severity of each kind of failure. The linear approach the court rejected is a sequence of threshold tests that discards options without comparing them, which can lead to adoption by default if each alternative has some risk. The holistic approach compares expected losses across all options at once, which is what minimising expected loss requires. The judgment in \emph{Re B-S} does not use the vocabulary of decision theory, but its method is a qualitative form of it.

\section{Private law: the welfare principle}

In disputes between parents about where a child should live and how much time she should spend with each, the child's welfare is the paramount consideration, guided by a checklist of factors in section 1 of the Children Act 1989. From 2014 the Act (section 1(2A)) included a presumption that the involvement of each parent in the child's life will further the child's welfare, unless the contrary is shown; the government announced in October 2025 that it would repeal the presumption. The presumption resembles a prior: in the absence of evidence of risk, the court starts from the view that both parents' involvement is good for the child. It is a prior with a gate, however: it is switched off, rather than updated, by some evidence of risk. Under section 1(6) the presumption does not apply to a parent where there is some evidence that the parent's involvement, in any form, would put the child at risk of harm, and the court must then assess the risk the parent poses. A review by the Ministry of Justice Harm Panel in 2020 raised the concern that the presumption could reinforce a culture favouring contact despite evidence of domestic abuse, and recommended that it be reviewed \citep{harmpanel2020}. A Bayesian treatment of such cases weighs the evidence of abuse, the predictive probability of future harm and the benefits of the relationship, and compares the expected welfare of the child under each arrangement.

\section{Risk assessment in child protection}

Child protection agencies have experimented with predictive risk models to help decide which reports of concern to investigate. A model developed for Allegheny County in Pennsylvania, which uses administrative data to score the risk that a child referred to the agency will later be removed from the home, has been studied extensively \citep{chouldechova2018,eubanks2018}. Such models raise the issues that Chapter~\ref{ch:lindleyai} takes up: calibration, disparate impact, the selective labels problem, since outcomes are observed only for children investigated, and the danger that the model reproduces patterns in past decisions rather than in children's actual risk. They also offer the possibility of more consistent screening decisions in a domain where workers face very high caseloads. The requirements for a Bayesian machine that Chapter~\ref{ch:lindleyai} will set out apply with special force when the decisions concern children.

\section*{Notes and Questions}
\addcontentsline{toc}{section}{Notes and Questions}
\begin{enumerate}[label=\arabic*.]
\item Explain the two level structure of the threshold criteria in section 31 as a pair of decision thresholds. Why is the threshold for future harm lower than that for past facts? Is the requirement that the prediction rest only on proved facts consistent with Bayesian updating?
\item An infant is injured while in the care of two parents and a babysitter. The court's posterior probabilities are $0.45$, $0.35$ and $0.20$. Who is in the pool of possible perpetrators? How should the court approach the question of returning the child to the first parent?
\item A skeletal finding occurs in thirty per cent of abused infants and one per cent of non abused infants with similar presentations. What is its likelihood ratio? With a prior of five per cent, what is the posterior probability of abuse?
\item Show how the holistic evaluation required by \emph{Re B-S} corresponds to comparing expected losses across options. Construct a numerical example with three options.
\item What would you require of a predictive risk model before using it to screen child protection referrals?
\end{enumerate}

\section*{Further reading}
\emph{Re H} \citep{reh1996}; \emph{Re S-B} \citep{resb2009}; \emph{Re B} (2013) \citep{reb2013}; \emph{Re B-S} \citep{rebs2013}.
